\section{Uniqueness via Canonical Representatives}
\label{sec:uniqueness}

\begin{figure}[t]
  \centering
  \begin{tikzpicture}[x=1cm,y=1cm,font=\small]
  \node[fig/node dark, text width=6.0cm] (assm) at (6.90,9.10)
    {unimodular $U$\\$U H_1 = H_2$\\$H_1, H_2$ in HNF};

  \node[fig/node, text width=3.6cm] (llu) at (2.90,7.10)
    {lower-left of $U$ vanishes\\[-0.1ex]{\footnotesize first-column equations}};
  \node[fig/node, text width=3.9cm] (llinv) at (10.90,7.10)
    {lower-left of $U^{-1}$ vanishes\\[-0.1ex]{\footnotesize same argument on $U^{-1} H_2 = H_1$}};

  \node[fig/node, text width=3.6cm] (lrU) at (2.90,5.15)
    {$\lowerRight(U)$ unimodular\\[-0.1ex]{\footnotesize restrict the inverse}};
  \node[fig/node, text width=3.6cm] (unit) at (10.90,5.15)
    {$U_{0,0}$ is a unit\\[-0.1ex]{\footnotesize read $(0,0)$ of $U U^{-1} = I$}};

  \node[fig/node dark, text width=5.3cm] (lreq) at (7.15,3.15)
    {$\lowerRight(H_1)=\lowerRight(H_2)$\\[-0.1ex]{\footnotesize induction hypothesis}};

  \node[fig/node, text width=3.2cm] (pivot) at (2.55,1.20)
    {pivot equality\\[-0.1ex]{\footnotesize normalized associates}};
  \node[fig/node, text width=3.4cm] (toprow) at (10.95,1.20)
    {top row equality\\[-0.1ex]{\footnotesize canonical remainders}};

  \node[fig/node plain, text width=2.4cm] (goal) at (6.90,-0.50)
    {$H_1 = H_2$};

  \coordinate (unitPivotBend) at (10.90,2.00);

  \draw[fig/arrow] (assm.south west) -- (llu.north);
  \draw[fig/arrow] (assm.south east) -- (llinv.north);
  \draw[fig/arrow] (llu.south) -- (lrU.north);
  \draw[fig/arrow] (llinv.south) -- (unit.north);
  \draw[fig/arrow] (lrU.south east) -- (lreq.north west);
  \draw[fig/arrow] (lreq.south) -- (pivot.north west);
  \draw[fig/arrow] (lreq.south east) -- (toprow.north west);
  \draw[fig/arrow] (unit.south) -- (unitPivotBend) -- (pivot.north east);
  \draw[fig/arrow] (pivot.south east) -- (goal.north west);
  \draw[fig/arrow] (toprow.south west) -- (goal.north east);
\end{tikzpicture}

  \caption{Layered dependency structure of the uniqueness proof.}
  \label{fig:uniqueness-graph}
\end{figure}

Section~\ref{sec:algorithm} established correctness: the executable procedure returns a
matrix in HNF. The deeper theorem is that this HNF is unique under left
equivalence once canonical remainder representatives are available.

\begin{theorem}
Let \(R\) be a Euclidean domain with \code{NormalizationMonoid},
decidable equality, and \code{CanonicalMod}. If \(H_1\) and \(H_2\) are matrices
in \code{IsHermiteNormalFormFin} and \(U\) is unimodular with
\(U \cdot H_1 = H_2\), then \(H_1 = H_2\).
\end{theorem}

The internal theorem is
\code{isHermiteNormalFormFin_unique_of_left_equiv}; the public theorem
\code{isHermiteNormalForm_unique_of_left_equiv} is obtained by reindexing.
The proof proceeds by induction on the recursive HNF witness for \(H_1\). The
empty-row and empty-column cases are immediate. The zero-column/pivot mismatch is
impossible because left multiplication preserves the property that the whole
first column is zero. The pivot/pivot case contains the real content and is
organized exactly as in Figure~\ref{fig:uniqueness-graph}.

\subsection{Lower-left vanishing and recursive descent}

Assume both \(H_1\) and \(H_2\) are in the pivot case. Because both matrices have
zeros below the first pivot, the equation \(U H_1 = H_2\) at entry \((i+1,0)\)
reduces to
\[
  U_{i+1,0} \cdot (H_1)_{0,0} = 0.
\]
Since \((H_1)_{0,0} \neq 0\), every \(U_{i+1,0}\) must be zero. This is the
lemma \code{lowerLeft_zero_of_mul_eq_pivot}. Once the lower-left block of
\(U\) vanishes, multiplication respects the block decomposition:
\[
  \lowerRight(U H_1) = \lowerRight(U) \cdot \lowerRight(H_1),
\]
formalized by \code{lowerRight_mul_of_lowerLeftZero}.

The same lower-left-zero fact also lets the proof transport unimodularity to the
lower-right block. The lemma \code{lowerRight_unimodular_of_lowerLeftZero}
builds a right inverse for \(\lowerRight(U)\) by restricting the inverse of
\(U\). At this point the induction hypothesis applies to the lower-right blocks,
yielding \(\lowerRight(H_1) = \lowerRight(H_2)\).

\subsection{Locking the pivot}

The next step is to identify the top-left pivot. Apply the same
lower-left-vanishing argument to \(U^{-1} H_2 = H_1\); this shows that the
lower-left block of \(U^{-1}\) is also zero. Evaluating
\code{Matrix.mul_nonsing_inv U hU} and
\code{Matrix.nonsing_inv_mul U hU} at the \((0,0)\) entry therefore gives
\[
  U_{0,0} \cdot (U^{-1})_{0,0} = 1,
  \qquad
  (U^{-1})_{0,0} \cdot U_{0,0} = 1.
\]
Hence \(U_{0,0}\) is a unit.

Now inspect the top-left entry of \(U H_1 = H_2\). Because the lower rows of
\(H_1\) are zero in the first column,
\[
  (H_2)_{0,0} = U_{0,0} \cdot (H_1)_{0,0}.
\]
Thus the two pivots are associated. Both pivots are normalized by the HNF
predicate, so \code{Associated.eq_of_normalized} yields
\((H_1)_{0,0} = (H_2)_{0,0}\). Substituting this equality back into the display
above gives
\[
  (U_{0,0} - 1) \cdot (H_1)_{0,0} = 0.
\]
Since the pivot is nonzero, \(U_{0,0} = 1\). This is the algebraic reason that
the top row cannot be rescaled nontrivially.

\subsection{Locking the top row by canonical remainders}

After the pivot is fixed and the lower-right blocks are identified, the first row
of the equation \(U H_1 = H_2\) becomes
\[
  H_2[0,\ast] = H_1[0,\ast] + \sum_i U_{0,i+1} \, H_1[i+1,\ast].
\]
All lower rows on the right-hand side are already known to coincide with the
corresponding lower rows of \(H_2\). The remaining problem is therefore to show
that every coefficient \(U_{0,i+1}\) is zero, so that the top rows are equal.

This is where \code{CanonicalMod} becomes indispensable. The proof
\code{topRow_eq_of_eq_add_lowerRows} assumes, for contradiction, that some
coefficient is nonzero and chooses the least active lower row whose trailing part
has a pivot. Let that pivot occur at column \(p+1\). Minimality and the HNF zero
structure imply that all other summands vanish in that column, so the displayed
equation reduces there to
\[
  (H_2)_{0,p+1} = (H_1)_{0,p+1} + c \cdot (H_1)_{i+1,p+1}.
\]
Because both matrices are already in HNF, the entries
\((H_1)_{0,p+1}\) and \((H_2)_{0,p+1}\) are already reduced modulo the common
pivot \((H_1)_{i+1,p+1}\). Their difference is divisible by that pivot, so
\code{CanonicalMod.reduced_eq_of_dvd_sub} forces the two entries to be equal.
The supposedly nonzero coefficient \(c\) is therefore annihilated by a nonzero
pivot, contradiction.

Two points are worth isolating.

\begin{itemize}
  \item The single-step reduction proof in the algorithm used only
    \code{CanonicalMod.mod_mod}. Uniqueness needs the strictly stronger field
    \code{reduced_eq_of_dvd_sub}.
  \item \code{NormalizationMonoid} is still essential, but for a different
    reason. It normalizes pivots up to associates; it does not control exact
    remainder representatives. The two abstractions are complementary, not
    redundant.
\end{itemize}

This separation is the main conceptual payoff of the mixin design. The Euclidean
algorithm computes with whatever remainder operation the domain provides, while
the uniqueness theorem imports only the exact canonicality principles it needs to
state equality rather than equality up to associates.
